\section{Conclusion}
The Morlet CNN--Conformer produced the strongest aggregate accuracy/F1 profile among the four architectures evaluated under the fixed subject-separated MIT-BIH protocol. That result, however, came with a clear boundary: recognition remained concentrated in Normal and VEB, while SVEB and Fusion performance stayed near zero and was not rescued by a simple post-hoc threshold change. The study is therefore most useful as an auditable benchmark rather than as evidence of clinical readiness. Its fixed record separation, archived predictions, comparator runs, and class-wise diagnostics provide a reproducible basis for testing methods that specifically target rare-rhythm recognition. Progress toward clinician-supervised use will require better minority-class performance, independently selected operating points, external validation, and prospective assessment of how model outputs affect real review decisions.
